\documentclass[../main.tex]{subfiles}

\begin{document}

\chapter{System of Linear Equations}

One of the most important application of matrices would be in
linear systems of equations. With the ideas on matrices, we can
open our third eye when viewing system of linear equations.
This note is dedicated for matrices in linear equations,
Gaussian Elimination, and LU Decompositions.

\section{Basics and Intuitions}

As always, let's get started with fundamental definitions, concepts,
and intuitions that will be used throughout.

First, I am pretty sure most of you are aware of what system of
equation is. However, because it is always nice to formally state
our understanding, here is the definition of linear system of
equation.

\begin{keydef}
A \textbf{system of linear equations} is composed of linear equations,
or equations whose highest degree is $1$. A system of $m$ equations
and $n$ variables adhere to the following forms where $a_{ij}$ and
$b_{i}$ are known quantities.
\begin{align*}
    a_{11}x_1 + a_{12}x_2 + a_{13}x_3 + \cdots + a_{1n}x_n &= b_1 \\
    a_{21}x_1 + a_{22}x_2 + a_{23}x_3 + \cdots + a_{2n}x_n &= b_2 \\
    a_{31}x_1 + a_{32}x_2 + a_{33}x_3 + \cdots + a_{3n}x_n &= b_3 \\
    \vdots\qquad\qquad& \\
    a_{m1}x_1 + a_{m2}x_2 + a_{m3}x_3 + \cdots + a_{mn}x_n &= b_m
\end{align*}
\end{keydef}

Below is an example of a system of linear equations.
\begin{align*}
    x + 2y + 3z &= 4 \\
    -x + 3y + 2z &= 5
\end{align*}
Notice that all equations above are linear, meaning, their degrees
are $1$. Now, what does it mean to find a solution to such systems?

\begin{keydef}
A \textbf{solution to a system of linear equations} is a set of
scalar values for unknown variables such that all linear equations
are satisfied when substituted.
\end{keydef}

For instance, $(2, 1)$ is the solution to the following system of
linear equations.
\begin{align*}
    x - 2y &= 0 \\
    x + y  &= 3
\end{align*}
To check it, you can simply substitute to see if the equations are
satisfied.

Now, if you recall from the previous note where we discussed matrix
multiplication, system of linear equations can be written as a
matrix multiplication.

\begin{keydef}
A system of linear equation
\begin{align*}
    a_{11}x_1 + a_{12}x_2 + a_{13}x_3 + \cdots + a_{1n}x_n &= b_1 \\
    a_{21}x_1 + a_{22}x_2 + a_{23}x_3 + \cdots + a_{2n}x_n &= b_2 \\
    a_{31}x_1 + a_{32}x_2 + a_{33}x_3 + \cdots + a_{3n}x_n &= b_3 \\
    \vdots\qquad\qquad& \\
    a_{m1}x_1 + a_{m2}x_2 + a_{m3}x_3 + \cdots + a_{mn}x_n &= b_m
\end{align*}
can be written as $Ax = b$ where $A$, $x$, and $b$ are matrices
defined as the following.
\[
A =
\begin{bmatrix}
    a_{11} & a_{12} & a_{13} & \cdots & a_{1n} \\
    a_{21} & a_{22} & a_{23} & \cdots & a_{2n} \\
    a_{31} & a_{32} & a_{33} & \cdots & a_{3n} \\
    \vdots & \vdots & \vdots & \ddots & \vdots \\
    a_{m1} & a_{m2} & a_{m3} & \cdots & a_{mn} \\
\end{bmatrix},
\quad
x =
\begin{bmatrix}
    x_1 \\
    x_2 \\
    x_3 \\
    \vdots \\
    x_n
\end{bmatrix},
\quad
b
\begin{bmatrix}
    b_1 \\
    b_2 \\
    b_3 \\
    \vdots \\
    b_m
\end{bmatrix}
\]
\end{keydef}

When you actually perform matrix multiplication, you can see that
$Ax = b$ is identical to the system of linear equations. Now that
we have basic understanding of the relationship between matrices
and linear system of equations, let's take a look at a theorem and
an important corollary.

\begin{theorem}{Different Solutions to a System}
{Different Solutions to a System}
\label{Different Solutions to a System}
For two distinct solutions $x_1$ and $x_2$ of $Ax = b$,
$\alpha x_1 + \beta x_2$ is also a solution to the linear system
given that $\alpha, \beta \in \mathbb{R}$ and $\alpha + \beta = 1$.
\end{theorem}
\begin{proof}
From Theorem \ref{Properties of Matrix Multiplication}, it was shown
that $\lambda (AB) = A (\lambda B)$ for scalar $\lambda$ and matrices
$A$ and $B$ that can be multiplied. Substituting $\alpha x_1 +
\beta x_2$, the following system is obtained.
\[
A (\alpha x_1 + \beta x_2)
= A(\alpha x_1) + A(\beta x_2)
= \alpha (Ax_1) + \beta (Ax_2)
= \alpha b + \beta b
= b
\]
Therefore, $(\alpha x_1 + \beta x_2)$ is another solution to the
system.
\end{proof}

From this theorem, we can establish an important corollary about the
number of solutions to a system of linear equations.

\begin{corollary}{Number of Solutions to a Linear System}
{Number of Solutions to a Linear System}
\label{Number of Solutions to a Linear System}
A system of linear equations can retain $0$, $1$, or infinitely many
solutions.
\end{corollary}
\begin{proof}
First, it is self-evident that a system can retain $0$ or $1$
solution. Therefore, the fact that there exists infinite solutions
if there are more than $1$ solution must be proven.

Let $f(\alpha) = \alpha x_1 + (1 - \alpha) x_2$ be a function that
returns possible solutions where $x_1 \neq x_2$ are fixed values.
It suffices to show that the function $f$ is injective as there exists
infinitely many values of $\alpha$. Consider the following equations.
\begin{align*}
    f(\alpha)
    &= \alpha x_1 + (1 - \alpha) x_2 
     = \alpha (x_1 - x_2) + x_2 \\
    f(\alpha')
    &= \alpha' x_1 + (1 - \alpha') x_2 
     = \alpha' (x_1 - x_2) + x_2
\end{align*}
If $f(\alpha) = f(\alpha')$, then $\alpha (x_1 - x_2) + x_2 =
\alpha' (x_1 - x_2) + x_2$ and $\alpha = \alpha'$, proving the
injectivity. Therefore, a system of linear equations can only have
$0$, $1$, or infinitely many solutions.
\end{proof}

In middle school, you probably saw a visual understanding of the
corollary above. Let's do the same, except in 3-dimensional space.
Below are two examples of when there exist no solution. Note that
not all three planes intersect at a single point.

\begin{minipage}{0.48\textwidth}
\tdplotsetmaincoords{80}{100}
\begin{tikzpicture}[
        tdplot_main_coords,
        axis/.style={->,thick},
        scale=0.55
    ]
    \pgfmathsetmacro{\a}{1.13397459622}
    \pgfmathsetmacro{\b}{6.33012701892}

    \foreach \x in {0,1,...,9}
        \foreach \y in {0,1,...,9} {
			\draw[gray,very thin]
                (\x,0,0) -- (\x,9.5,0);
			\draw[gray,very thin] (0,\y,0) -- (9.5,\y,0);
		}
    \foreach \y in {0,1,...,9}
        \foreach \z in {0,1,...,9} {
			\draw[gray,very thin]
                (0,\y,0) -- (0,\y,9.5);
			\draw[gray,very thin] (0,0,\z) -- (0,9.5,\z);
		}
    \foreach \z in {0,1,...,9}
        \foreach \x in {0,1,...,9} {
			\draw[gray,very thin]
                (0,0,\z) -- (9.5,0,\z);
			\draw[gray,very thin] (\x,0,0) -- (\x,0,9.5);
		}

	\draw[axis] (-2,0,0) -- (10,0,0) node[left] {$x$};
	\draw[axis] (0,-1,0) -- (0,10,0) node[right] {$y$};
	\draw[axis] (0,0,-1) -- (0,0,10) node[left] {$z$};

	\draw[opacity=.5,very thick,fill=Red]
        (2,2,2) -- (8,2,2) -- (8,8,2) -- (2,8,2) -- cycle;
	\draw[opacity=.5,very thick,fill=Blue]
        (2,2.5,\a) -- (8,2.5,\a) -- (8,5.5,\b) -- (2,5.5,\b) -- cycle;
	\draw[opacity=.5,very thick,fill=yellow]
        (2,7.5,\a) -- (8,7.5,\a) -- (8,4.5,\b) -- (2,4.5,\b) -- cycle;
    \draw[very thick, dashed] (2,3,2) -- (8,3,2);
    \draw[very thick, dashed] (2,7,2) -- (8,7,2);
    \draw[very thick, dashed] (2,5,5.46410161514) --
        (8,5,5.46410161514);
\end{tikzpicture}
\end{minipage}
\hfill
\begin{minipage}{0.48\textwidth}
\tdplotsetmaincoords{80}{100}
\begin{tikzpicture}[
        tdplot_main_coords,
        axis/.style={->,thick},
        scale=0.55
    ]
    \foreach \x in {0,1,...,9}
        \foreach \y in {0,1,...,9} {
			\draw[gray,very thin]
                (\x,0,0) -- (\x,9.5,0);
			\draw[gray,very thin] (0,\y,0) -- (9.5,\y,0);
		}
    \foreach \y in {0,1,...,9}
        \foreach \z in {0,1,...,9} {
			\draw[gray,very thin]
                (0,\y,0) -- (0,\y,9.5);
			\draw[gray,very thin] (0,0,\z) -- (0,9.5,\z);
		}
    \foreach \z in {0,1,...,9}
        \foreach \x in {0,1,...,9} {
			\draw[gray,very thin]
                (0,0,\z) -- (9.5,0,\z);
			\draw[gray,very thin] (\x,0,0) -- (\x,0,9.5);
		}

	\draw[axis] (-2,0,0) -- (10,0,0) node[left] {$x$};
	\draw[axis] (0,-1,0) -- (0,10,0) node[right] {$y$};
	\draw[axis] (0,0,-1) -- (0,0,10) node[left] {$z$};

	\draw[opacity=.5,very thick,fill=Red]
        (2,2,2) -- (8,2,2) -- (8,8,2) -- (2,8,2) -- cycle;
	\draw[opacity=.5,very thick,fill=Blue]
        (2,2,4) -- (8,2,4) -- (8,8,4) -- (2,8,4) -- cycle;
	\draw[opacity=.5,very thick,fill=Yellow]
        (2,2,6) -- (8,2,6) -- (8,8,6) -- (2,8,6) -- cycle;
\end{tikzpicture}
\end{minipage}

Continuing, the left diagram represents the case where there
is a single solution while the right diagram shows the system
with infinite solutions.

\begin{minipage}{0.48\textwidth}
\tdplotsetmaincoords{80}{100}
\begin{tikzpicture}[
        tdplot_main_coords,
        axis/.style={->,thick},
        scale=0.55
    ]
    \pgfmathsetmacro{\a}{6.24264068712}

    \foreach \x in {0,1,...,9}
        \foreach \y in {0,1,...,9} {
			\draw[gray,very thin]
                (\x,0,0) -- (\x,9.5,0);
			\draw[gray,very thin] (0,\y,0) -- (9.5,\y,0);
		}
    \foreach \y in {0,1,...,9}
        \foreach \z in {0,1,...,9} {
			\draw[gray,very thin]
                (0,\y,0) -- (0,\y,9.5);
			\draw[gray,very thin] (0,0,\z) -- (0,9.5,\z);
		}
    \foreach \z in {0,1,...,9}
        \foreach \x in {0,1,...,9} {
			\draw[gray,very thin]
                (0,0,\z) -- (9.5,0,\z);
			\draw[gray,very thin] (\x,0,0) -- (\x,0,9.5);
		}

	\draw[axis] (-2,0,0) -- (10,0,0) node[left] {$x$};
	\draw[axis] (0,-1,0) -- (0,10,0) node[right] {$y$};
	\draw[axis] (0,0,-1) -- (0,0,10) node[left] {$z$};

	\draw[opacity=.5,very thick,fill=Red]
        (2,2,2) -- (2,2,8) -- (\a,\a,8) -- (\a,\a,2) -- cycle;
	\draw[opacity=.5,very thick,fill=Blue]
        (\a,2,2) -- (2,\a,2) -- (2,\a,8) -- (\a,2,8) -- cycle;
	\draw[opacity=.5,very thick,fill=Yellow]
        (2,2,5) -- (2,\a,5) -- (\a,\a,5) -- (\a,2,5) -- cycle;
    \draw[very thick, dashed] (4.12132034356,4.12132034356,2) --
        (4.12132034356,4.12132034356,8);
    \draw[fill] (4.12132034356,4.12132034356,5) circle (3pt);
\end{tikzpicture}
\end{minipage}
\hfill
\begin{minipage}{0.48\textwidth}
\tdplotsetmaincoords{80}{100}
\begin{tikzpicture}[
        tdplot_main_coords,
        axis/.style={->,thick},
        scale=0.55
    ]
    \pgfmathsetmacro{\a}{1.40192378865}
    \pgfmathsetmacro{\b}{6.59807621135}

    \foreach \x in {0,1,...,9}
        \foreach \y in {0,1,...,9} {
			\draw[gray,very thin]
                (\x,0,0) -- (\x,9.5,0);
			\draw[gray,very thin] (0,\y,0) -- (9.5,\y,0);
		}
    \foreach \y in {0,1,...,9}
        \foreach \z in {0,1,...,9} {
			\draw[gray,very thin]
                (0,\y,0) -- (0,\y,9.5);
			\draw[gray,very thin] (0,0,\z) -- (0,9.5,\z);
		}
    \foreach \z in {0,1,...,9}
        \foreach \x in {0,1,...,9} {
			\draw[gray,very thin]
                (0,0,\z) -- (9.5,0,\z);
			\draw[gray,very thin] (\x,0,0) -- (\x,0,9.5);
		}

	\draw[axis] (-2,0,0) -- (10,0,0) node[left] {$x$};
	\draw[axis] (0,-1,0) -- (0,10,0) node[right] {$y$};
	\draw[axis] (0,0,-1) -- (0,0,10) node[left] {$z$};

	\draw[opacity=.5,very thick,fill=Red]
        (2,2,4) -- (8,2,4) -- (8,8,4) -- (2,8,4) -- cycle;
	\draw[opacity=.5,very thick,fill=Blue]
        (2,3.5,\a) -- (8,3.5,\a) -- (8,6.5,\b) -- (2,6.5,\b) -- cycle;
	\draw[opacity=.5,very thick,fill=Yellow]
        (2,3.5,\b) -- (8,3.5,\b) -- (8,6.5,\a) -- (2,6.5,\a) -- cycle;
    \draw[very thick] (2,5,4) -- (8,5,4);
\end{tikzpicture}
\end{minipage}

Wow, not going to lie, it took some time to draw the diagrams above.
Anyhow, with that in mind, let's discuss other fundamental terms
and concepts before we move on to the next section of the note.

\begin{keydef}
A linear system is \textbf{consistent} if it retains at least one
solution, whereas a system is \textbf{inconsistent} if there is no
solution.
\end{keydef}

We could also define a system based on the column matrix $b$.

\begin{keydef}
A linear system is \textbf{homogeneous} if the column matrix $b$
in the linear system $Ax = b$ is a zero matrix. If the matrix is not
a zero matrix, the system is \textbf{nonhomogeneous}, or
\textbf{inhomogeneous}.
\end{keydef}

Now from the definitions, we can see that all homogeneous systems
are consistent. As you have guessed, we have a special name for
a special solution in such cases.

\begin{keydef}
\textbf{Trivial solution} is a solution where all the entries for
$x$ in $Ax = b$ are zero.
\end{keydef}

With the fundamental terms in mind, let's take a look at how matrices
can be used to solve complex linear system with Gaussian Elimination.

\section{Gaussian Elimination}

When you solve a linear system, you probably multiplied constants
and add or subtract equations to cancel variables. That is exactly
what we will be doing in Gaussian Elimination, but with matrices.
First, define how we solve a linear system in a more rigorous way.

\begin{keydef}
For a system of linear equations, performing following operations
will result in a linear system with identical solution.
\begin{enumerate}[noitemsep]
\item
Change the order of linear equations in which they are listed
\item
Multiply a nonzero constant, or nonzero scalars, to an equation
\item
Add equations from the system to form a new equation
\end{enumerate}
\end{keydef}

Now before we connect this idea with matrices, let's define a matrix
that we will use with the property of a linear system above.

\begin{keydef}
An \textbf{augmented matrix} for a linear system $Ax = b$ is a
partitioned matrix in the form $[A \mid b]$.
\end{keydef}

Consider the following system for an example.
\begin{align*}
    x + 2y + 3z &= 4 \\
    2x + 3y + 4z &= 5 \\
    3x + 4y + 5z &= 6
\end{align*}
In the linear system above, the augmented matrix will be the following
matrix.
\[
\begin{bNiceArray}{ccc|c}
    1 & 2 & 3 & 4 \\
    2 & 3 & 4 & 5 \\
    3 & 4 & 5 & 6
\end{bNiceArray}
\]
Now that we have augmented matrix in mind, let's take a look at
how the three properties that we discussed applies with matrices.

\begin{keydef}
For an augmented matrix for a linear system, performing the following
operations will result in a linear system with identical solution as
the initial system.
\begin{enumerate}[noitemsep]
\item
Rearrangement of the rows in the augmented matrix
\item
Multiply a nonzero scalar to a row in the augmented matrix
\item
Add a row multiplied by a scalar with another to replace a row with
new row
\end{enumerate}
Such operations are referred to as \textbf{elementary row operations}.
\end{keydef}

If you compare the operations one by one, you can notice that the
two set of properties refer to the same property. Now with the
set of property of augmented matrices in mind, here are the steps
for the Gaussian Elimination algorithm.

\begin{theorem}{Gaussian Elimination}{Gaussian Elimination}
\label{Gaussian Elimination}
Execute the following steps in order for Gaussian Elimination.
\begin{enumerate}[noitemsep]
\item
Construct the augmented matrix.
\item
Utilize the elementary row operations to transform the augmented
matrix into REF.
\item
Convert the augmented matrix in REF into equations of a transformed
linear system, referred to as the \textbf{derived set}.
\item
Back substitute to solve the system.
\end{enumerate}
\end{theorem}

Let's take a look at few examples.

\begin{problem}{Gaussian Elimination Problem 1}
{Gaussian Elimination Problem 1}
\label{Gaussian Elimination Problem 1}
Find solution(s) to the following system of linear equations if any.
\begin{align*}
    x + 2y + z &= 1 \\
    3x - y + z &= 6 \\
    x + y - z &= -2
\end{align*}
\end{problem}
\begin{solution}
First, the augmented matrix can be found from the system. Consider
the matrix below.
\[
\begin{bNiceArray}{ccc|c}
    1 & 2  & 1  & 1 \\
    3 & -1 & 1  & 6 \\
    1 & 1  & -1 & -2
\end{bNiceArray}
\]
Using the elementary row operations, the augmented matrix can be
transformed as the following.
\begin{align*}
\begin{bNiceArray}{ccc|c}
    1 & 2  & 1  & 1 \\
    3 & -1 & 1  & 6 \\
    1 & 1  & -1 & -2
\end{bNiceArray}
&\rightarrow
\begin{bNiceArray}{ccc|c}
    1 & 2  & 1  & 1 \\
    0 & -4 & 4  & 12 \\
    1 & 1  & -1 & -2
\end{bNiceArray}
\rightarrow
\begin{bNiceArray}{ccc|c}
    1 & 2  & 1  & 1 \\
    0 & -4 & 4  & 12 \\
    0 & -1 & -2 & -3
\end{bNiceArray} \\
&\rightarrow
\begin{bNiceArray}{ccc|c}
    1 & 2  & 1  & 1 \\
    0 & -1 & 1  & 3 \\
    0 & 1  & 2  & 3
\end{bNiceArray}
\rightarrow
\begin{bNiceArray}{ccc|c}
    1 & 2  & 1  & 1 \\
    0 & -1 & 1  & 3 \\
    0 & 0  & 3  & 6
\end{bNiceArray}
\rightarrow
\begin{bNiceArray}{ccc|c}
    1 & 2  & 1  & 1 \\
    0 & 1  & -1 & -3 \\
    0 & 0  & 1  & 2
\end{bNiceArray}
\end{align*}
Because the augmented matrix is now in REF, the transformed
linear system can be rewritten.
\begin{align*}
    x + 2y + z &= 1 \\
    y - z &= -3 \\
    z &= 2
\end{align*}
Back substituting, it is evident that $z = 2$ and $y = -1$.
Moreover, $x = 1 + 2 - 2 = 1$. Therefore, the triple $(1, -1, 2)$
is the solution to the linear system.
\end{solution}

Try substituting to the original system! You can notice that the
solution is valid. Below is the second example of using
Gaussian Elimination.

\begin{problem}{Gaussian Elimination Problem 2}
{Gaussian Elimination Problem 2}
\label{Gaussian Elimination Problem 2}
Find solution(s) to the following system of linear equations if any.
\begin{align*}
    x + 2y + 3z &= 4 \\
    y + 3z &= 5 \\
    x + y &= 1
\end{align*}
\end{problem}
\begin{solution}
First and foremost, the augmented matrix can be found.
\[
\begin{bNiceArray}{ccc|c}
    1 & 2 & 3 & 4 \\
    0 & 1 & 3 & 5 \\
    1 & 1 & 0 & 1
\end{bNiceArray}
\rightarrow
\begin{bNiceArray}{ccc|c}
    1 & 2  & 3  & 4 \\
    0 & 1  & 3  & 5 \\
    0 & -1 & -3 & -3
\end{bNiceArray}
\rightarrow
\begin{bNiceArray}{ccc|c}
    1 & 2 & 3 & 4 \\
    0 & 1 & 3 & 5 \\
    0 & 0 & 0 & 2
\end{bNiceArray}
\]
Writing the transformed augmented matrix in to a system of linear
equations, the following system is obtained.
\begin{align*}
    x + 2y + 3z &= 4 \\
    y + 3z &= 5 \\
    0 &= 2
\end{align*}
Notice that $0 = 2$ can never be true. Therefore, no solution exists.
\end{solution}

We could double check out answer without using matrices. Notice that
$3z - x = 4$ and $x = 3z - 4$. Similarly, $y = 5 - 3z$. Substituting
to the first equation, $3z - 4 + 10 - 6z + 3z = 4$, or $6 = 4$,
which is not true. Therefore, our results align! Below is the third
example of using the algorithm.

\begin{problem}{Gaussian Elimination Problem 3}
{Gaussian Elimination Problem 3}
\label{Gaussian Elimination Problem 3}
Find solution(s) to the following system of linear equations if any.
\begin{align*}
    x - y - 2z &= 3 \\
    2x + y - 7z &= 3 \\
    x + y - 4z &= 1
\end{align*}
\end{problem}
\begin{solution}
Consider the following augmented matrix.
\[
\begin{bNiceArray}{ccc|c}
    1 & -1 & -2 & 3 \\
    2 & 1  & -7 & 3 \\
    1 & 1  & -4 & 1
\end{bNiceArray}
\]
Utilizing Gaussian Elimination, the augmented matrix can be
transformed as following.
\begin{align*}
\begin{bNiceArray}{ccc|c}
    1 & -1 & -2 & 3 \\
    2 & 1  & -7 & 3 \\
    1 & 1  & -4 & 1
\end{bNiceArray}
&\rightarrow
\begin{bNiceArray}{ccc|c}
    1 & -1 & -2 & 3 \\
    2 & 1  & -7 & 3 \\
    0 & 2  & -2 & -2
\end{bNiceArray}
\rightarrow
\begin{bNiceArray}{ccc|c}
    1 & -1 & -2 & 3 \\
    0 & 1  & -1 & -1 \\
    2 & 1  & -7 & 3
\end{bNiceArray} \\
&\rightarrow
\begin{bNiceArray}{ccc|c}
    1 & -1 & -2 & 3 \\
    0 & 1  & -1 & -1 \\
    0 & 3  & -3 & -3
\end{bNiceArray}
\rightarrow
\begin{bNiceArray}{ccc|c}
    1 & -1 & -2 & 3 \\
    0 & 1  & -1 & -1 \\
    0 & 0  & 0  & 0
\end{bNiceArray}
\end{align*}
Rewriting the system, the following equations are obtained.
\begin{align*}
    x - y - 2z &= 3 \\
    y - z &= -1
\end{align*}
Notice that $y = z - 1$ and $x = y + 2z + 3 = z - 1 + 2z + 3 =
3z + 2$. Because there exists infinitely many $z$, the system
has infinitely many solutions.
\end{solution}

When writing such solutions, it could be written as the following.
\[
x =
\begin{bmatrix}
    3z + 2 \\
    z - 1 \\
    z
\end{bmatrix}
\]























\end{document}



